\section{Material related to learning-theoretic analysis}
\subsection{Proof of Theorem~\ref{thm:effective-learning}}\label{subapp:proof-thm-effective-learning}
Before we provide the proof of this theorem, we first describe in Section~\ref{subsubsection:useful-results} how a state $U^T(\rvx)_t =: \rho_t$ depends on previous states in Claim~\ref{cl:closed-form-recursion}, followed by a set-injectivity result in Proposition~\ref{prop:set-injectivity} that we leverage to establish the proof of Theorem~\ref{thm:effective-learning}.


\subsubsection{Useful results}\label{subsubsection:useful-results}
\begin{claim}[Closed form reservoir recursion]\label{cl:closed-form-recursion}
Let $\rho_+ = \pdensity$ and define the input-dependent channel $\gJ_{x_t}(\rho) := \gJ(\rho, x_t)$ for elements $x_t$ of a window $\rvx = (x_{\tau-w+1}, \ldots, x_{\tau})$. Then the CEQC recursion of~\cref{eq:CEQC} can be written as
\begin{equation*}
    \rho_t = \lambda \gJ_{x_t}(\rho_{t-1}) + (1-\lambda) \rho_+.
\end{equation*}
Define for $k \geq 1$ the backward composition operator
\begin{equation*}
    \gJarr{k}_{t}(\rho | \rvx) := \pth{\gJ_{x_t} \circ \gJ_{x_{t-1}} \circ \cdots \circ \gJ_{x_{t-k+1}}} (\rho),
\end{equation*}
and for $k = 0$, $\gJarr{k}_{t}(\rho|\rvx) := \rho$. Then
the reservoir's state after consuming inputs $(x_{\tau - w + 1}, \dots, x_t)$ becomes
\begin{equation}\label{eq:recursion}
    \rho_t = U^T(\rvx)_t = \lambda^m \gJarr{m}_{t}(\rho_{\tau-w} | \rvx) + (1-\lambda) \sum_{k=0}^{m-1} \lambda^k \gJarr{k}_{t}(\rho_+ | \rvx),
\end{equation}
with $m$ being the number of applications of the evolution $T$: $m = t - (\tau - w) \in \{1, \ldots, w\}$.
\end{claim}
\begin{proof}
    By induction on $m$. The case $m=1$ is~\cref{eq:CEQC}. Assuming the formula holds at time $t-1$, substitute it into $\rho_t = \lambda \gJ_{x_t}(\rho_{t-1}) + (1-\lambda) \rho_+$ and use the definition of $\gJarr{k}_{t}$ to obtain~\cref{eq:recursion}.
\end{proof}

\begin{proposition}[Set-injectivity]\label{prop:set-injectivity}
    Let $\gD = \{\rvx_1, \ldots, \rvx_N\}$ be a dataset of windows. Suppose our quantum featurizer (initialized at $\rho_{-w} := \pdensity$) contains $R \geq 1$ sub-reservoirs that follow the design described in Section~\ref{subsec:QR-embedding}. Let our readout scheme follow the kernel-based one described in Section~\ref{subsec:K-read}. Let $\varepsilon_{\pr}, \delta_{\pr}  \in (0,1)$ be, respectively, a tolerance and a failure probability choice. Then, with a number of qubits as low as $n = \Omega\pth{\varepsilon_{\pr}^{-2} \log\pth{\tfrac{w N}{\delta_{\pr}}}}$, we have $m \circ H^T(\rvx) \neq m \circ H^T(\rvx')$ for any pair of windows $\rvx \neq \rvx' \in \gD$ with probability at least $1 - \delta_{\pr}$.
\end{proposition}

\begin{proof}
    In this proof, we only treat the case $R = 1$, as the cases $R > 1$ follow directly.
    
    Let $\gP = \bigcup_{i=1}^N \rvx_i \subset \sR^d$ be the set that contains all points that appear inside any window, which has cardinality at most $|\gP| \leq wN$. As described in Section~\ref{subsubsec:projection}, the event $\mathrm{E_{JL}}$ that describes that, for all $u, v \in \gP$,
    \begin{equation*}
        (1-\varepsilon_{\pr}) \normeuc{u - v}^2 \leq \normeuc{\Pi u - \Pi v}^2 \leq (1+\varepsilon_{\pr}) \normeuc{u - v}^2,
    \end{equation*}
    has probability $\Pr(\mathrm{E_{JL}}) \geq 1- \delta_{\pr}$ when $\Pi$ follows a Gaussian Johnson--Lindenstrauss distribution as stated in~\cref{eq:JL-matrix} and $n = \Omega\pth{\varepsilon_{\pr}^{-2} \log\pth{\tfrac{w N}{\delta_{\pr}}}}$, where the probability is over $\Pi$.
    Therefore, on the event $\mathrm{E_{JL}}$, the map $u \mapsto \Pi u$ is injective on the set $\gP$. From now on, condition on $\mathrm{E_{JL}}$.

    Let $\rvx = (x_{\tau - w + 1}, \ldots, x_\tau)$ and $\rvx' = (x'_{\tau' - w + 1}, \ldots, x'_{\tau'})$ be two distinct windows from the dataset $\gD$. To alleviate clutter, we use the re-indexing by relative lag $t \in \{-w+1, \ldots, 0\}$:
    \[
        x_t := x_{\tau + t}, \qquad x'_t := x'_{\tau' + t}.
    \]
    Since $\rvx$ and $\rvx'$ are distinct, let $t_\star$ be the latest lag where they differ:
    \begin{equation*}
        t_\star := \max\{t\in\{-w+1,\ldots,0\}: x_t \neq x'_t\}.
    \end{equation*}
    Then for all $t>t_\star$, we have $x_t=x'_t$. On the event $\mathrm{E_{JL}}$, $\Pi$ is injective on $\gP$, hence $z_{t_\star}\neq z'_{t_\star}$, where $z_t:=\Pi x_t$ and $z'_t:=\Pi x'_t$. In particular, there exists at least one coordinate $j_\star$ such that $(z_{t_\star})_{j_\star} \neq (z'_{t_\star})_{j_\star}$. Additionally, injectivity of the nonlinearity $\theta(u) = \pi \tanh(u)$ preserves this difference, so that $\theta\big((z_{t_\star})_{j_\star}\big) \neq \theta\big((z'_{t_\star})_{j_\star}\big)$.
    
    Let $\vartheta$ collect all random angles in the Ising unitary as well as the parameter $\lambda$ in~\cref{eq:contractive-channel}. Define the ``collision'' condition for this window pair $(\rvx, \rvx')$ as
    \begin{equation*}
        \gB_{\rvx, \rvx'}  := \{\vartheta: \Phi_{\vartheta}(\rvx) = \Phi_{\vartheta}(\rvx')\}, \quad \Phi_{\vartheta}(\rvx) := m \circ H^{T_\vartheta}(\rvx).
    \end{equation*}
    All that is left in the proof is to show that for each fixed pair $\rvx \neq \rvx'$, the set $\gB_{\rvx, \rvx'}$ has Lebesgue measure zero. A finite union argument then yields that the union of $\gB_{\rvx, \rvx'}$ over all $\binom{N}{2}$ pairs is still a measure-zero set; hence, with probability $1$ over $\vartheta$ (since $\vartheta$ is absolutely continuous with respect to. the Lebesgue measure under i.i.d.\ uniform sampling), no collisions occur on $\gD$ on the event $\mathrm{E_{JL}}$.
    

    Consider the scalar function
    \begin{equation*}
        g_{\rvx, \rvx'}(\vartheta) := \normeuc{\Phi_{\vartheta}(\rvx) - \Phi_{\vartheta}(\rvx')}^2.
    \end{equation*}
    Each feature coordinate $\Trc{O\,H^{T_\vartheta}(\rvx)}$ is obtained by composing finitely many maps that depend real-analytically on $\vartheta$ (products of single-/two-qubit rotations) with affine linear CPTP maps, and then taking a trace against a fixed observable. Hence $\Phi_{\vartheta}(\rvx)$ and therefore $g_{\rvx,\rvx'}(\vartheta)=\|\Phi_{\vartheta}(\rvx)-\Phi_{\vartheta}(\rvx')\|_2^2$ are real-analytic in $\vartheta$. Therefore, if we show a parameter choice $\vartheta^\star$ such that $g_{\rvx, \rvx'}(\vartheta^\star) > 0$, it follows from real-analyticity of $g_{\rvx, \rvx'}$ that the set $\{\vartheta: g_{\rvx, \rvx'}(\vartheta) = 0\}$ has Lebesgue measure zero (as can be found in Proposition~0 of~\cite{mityagin2015zeroset}). We argue in what follows the existence of such a witness.
    
    Let $\vartheta^\star_\bullet = 0$. Then we have 
    \begin{equation*}
        V(x) = \pth{\bigotimes_{j=1}^n R_y(\theta(z_j))}.
    \end{equation*}
    As each injection qubit undergoes only $Y$-rotations, if the two windows differ at the latest time $t_\star$ and projected coordinate $j_\star$, then the injected angles differ by
    \begin{equation*}
        \Delta \theta = \underbrace{\theta\big((z_{t_\star})_{j_\star} \big)}_{\alpha_{t_\star}} - \underbrace{\theta\big((z'_{t_\star})_{j_\star} \big)}_{\alpha'_{t_\star}} \neq 0.
    \end{equation*}
    Since $t_\star$ is the latest lag at which the windows differ, we have $\alpha_t = \alpha'_t$ for all $t > t_\star$, hence the accumulated angle $\phi := \sum_{t_\star + 1}^0 \alpha_t = \sum_{t_\star + 1}^0 \alpha'_t$. Consider the $1$-local observables $X_{j_\star}$ and $Z_{j_\star}$ at the level of qubit $j_\star$ (these belong to the feature vector since the measured set of $k$-local observables $\gO$ contains all $1$-local Paulis as discussed in Section~\ref{subsubsec:efficient-msmt-scheme}). Denote the sub-vectors
    \begin{equation*}
        \mu(\rvx) := \pth{\Braket{X_{j_\star}}_{\rvx}, \Braket{Z_{j_\star}}_{\rvx}} \in \sR^2, \quad \mu(\rvx') := \pth{\Braket{X_{j_\star}}_{\rvx'}, \Braket{Z_{j_\star}}_{\rvx'}} \in \sR^2.
    \end{equation*}
    Since
    $$
    \normeuc{\mu(\rvx) - \mu(\rvx')} \neq 0 \implies \|\Phi_{\vartheta}(\rvx)-\Phi_{\vartheta}(\rvx')\|_2 \neq 0,
    $$
    showing a witness $\lambda_\star \in (0,1)$ for which $\mu(\rvx) - \mu(\rvx') \neq 0$ establishes the proof of the proposition. This is precisely what we show in what follows.

    Leveraging the CEQC recursion defined in Claim~\ref{cl:closed-form-recursion}, the output state $H^T(\rvx)$ at lag $0$ admits the expansion
    \begin{equation*}
        H^T(\rvx) = \lambda^w \gJarr{w}_{0}(\rho_{-w}|\rvx) + (1-\lambda) \sum_{k=0}^{w-1} \lambda^k \gJarr{k}_{0}(\rho_+|\rvx),
    \end{equation*}
    where $\gJarr{k}_{0}(\cdot|\rvx)$ applies the last $k$ encoding steps (from lags $-k+1$ up to $0$) starting from state $\rho$. Define $i_\star := -t_\star + 1$. Because $\rvx$ and $\rvx'$ coincide for all lags $t > t_\star$, it follows that for all $k \leq i_\star - 1$,
    \begin{equation*}
        \gJarr{k}_{0}(\rho_+|\rvx) = \gJarr{k}_{0}(\rho_+|\rvx').
    \end{equation*}
    Hence the final state difference $\Delta \rho_0 := H^T(\rvx) - H^T(\rvx')$ can be written as
    \begin{equation}\label{eq:diff-state}
    \begin{aligned}
        \Delta \rho_0 = \; &\lambda^w \pth{\gJarr{w}_{0}(\rho_{-w}|\rvx) - \gJarr{w}_{0}(\rho_{-w}|\rvx')}\\
        &+ (1-\lambda) \lambda^{i_\star} \pth{\gJarr{i_\star}_{0}(\rho_+|\rvx) - \gJarr{i_\star}_{0}(\rho_+|\rvx')}\\
        &+ (1-\lambda) \sum_{k=i_\star + 1}^{w-1} \lambda^k \pth{\gJarr{k}_{0}(\rho_+|\rvx) - \gJarr{k}_{0}(\rho_+|\rvx')}.
    \end{aligned}
    \end{equation}
    Applying the linear map $\sigma \mapsto (\Trc{X_{j_\star}\sigma},\,\Trc{Z_{j_\star}\sigma})$ to~\cref{eq:diff-state}
    and using linearity of the trace yields
    \begin{equation}\label{eq:delta-mu-decomp}
    \mu(\rvx)-\mu(\rvx')
    =
    \lambda^w \Delta\mu^{(w)}_{\rho_{-w}}
    +
    (1-\lambda)\lambda^{i_\star} \Delta\mu^{(i_\star)}_{+}
    +
    (1-\lambda)\sum_{k=i_\star+1}^{w-1}\lambda^{k} \Delta\mu^{(k)}_{+},
    \end{equation}
    where, for any $k\ge 0$ and any starting state $\rho$,
    \[
    \begin{aligned}
        \Delta\mu^{(k)}_{\rho}
    :=
    \Big(
    &\Trc{X_{j_\star} \gJarr{k}_{0}(\rho|\rvx)} -
    \Trc{X_{j_\star} \gJarr{k}_{0}(\rho|\rvx')}, \\
    &\Trc{Z_{j_\star} \gJarr{k}_{0}(\rho|\rvx)}
    -
    \Trc{Z_{j_\star} \gJarr{k}_{0}(\rho|\rvx')}
    \Big).
    \end{aligned}
    \]
    
    As for the witness $\vartheta_\bullet^\star=0$, we have $W=\mathbb I$ and $\gJ_{x_t}(\rho)=V(x_t)\rho V(x_t)^\dagger$ with $V(x_t)=\bigotimes_{j=1}^n R_y(\alpha_t^{(j)})$. On qubit $j_\star$, starting from $\rho_+$ we have
    \[
    \Big(\Trc{X_{j_\star} \gJarr{k}_{0}(\rho_+|\rvx)},
    \Trc{Z_{j_\star} \gJarr{k}_{0}(\rho_+|\rvx)}\Big)
    =
    (\cos(\phi+\alpha_{t_\star}), \sin(\phi+\alpha_{t_\star})),
    \]
    and similarly with $\alpha'_{t_\star}$ for $\rvx'$.
    Hence we have
    \begin{equation}\label{eq:lead-delta}
    \begin{aligned}
        \big\|\Delta\mu^{(i_\star)}_{+}\big\|_2
        \!&=\!
        \normeuc{(\cos(\phi\!+\!\alpha_{t_\star}),\sin(\phi\!+\!\alpha_{t_\star}))
        \!-\!(\cos(\phi+\alpha'_{t_\star}),\sin(\phi+\alpha'_{t_\star}))}\\
        &=
        2\big|\sin(\Delta\theta/2)\big|
        >0.
    \end{aligned}
    \end{equation}

    For any state $\sigma$, $|\Trc{X_{j_\star}\sigma}|\le 1$ and $|\Trc{Z_{j_\star}\sigma}|\le 1$;
    therefore $\|\Delta\mu^{(k)}_{\rho}\|_2\le 2\sqrt{2}$ for all $k$ and all $\rho$.
    Using~\cref{eq:delta-mu-decomp} and the triangle inequality gives, for $i_\star \le w-1$,
    \[
    \begin{aligned}
        \|\mu(\rvx)-\mu(\rvx')\|_2
        \ge\;&(1-\lambda)\lambda^{i_\star}\cdot 2|\sin(\Delta\theta/2)|\\
        &-
        2\sqrt{2}\lambda^w
        -
        2\sqrt{2}(1-\lambda)\sum_{k=i_\star+1}^{w-1}\lambda^{k}.
    \end{aligned}
    \]
    Since $(1-\lambda)\sum_{k=i_\star+1}^{w-1}\lambda^{k}\le \lambda^{i_\star+1} - \lambda^w$, we obtain
    \[
    \begin{aligned}
    \normeuc{\mu(\rvx)-\mu(\rvx')}
    \ge 2(1-\lambda)\lambda^{i_\star}|\sin(\Delta\theta/2)| - 2\sqrt{2}\lambda^{i_\star+1}.
    \end{aligned}
    \]
    Choosing $\lambda_\star \in (0,1/2]$ small enough so that
    \[
    \lambda_\star < \frac{|\sin(\Delta\theta/2)|}{\sqrt{2} + |\sin(\Delta\theta/2)|},
    \]
    yields $\mu(\rvx)\neq\mu(\rvx')$. (If $i_\star=w$, i.e.\ $t_\star=-w+1$, the same computation applies with the leading contribution coming from the $\lambda^w \Delta\mu^{(w)}_{\rho_{-w}}$ term.) This concludes the proof.
    
\end{proof}

\subsubsection{Proof of Theorem~\ref{thm:effective-learning}}\label{subsubsection:proof-thm-effective-learning}
Recall the setup of the K-read training in Section~\ref{subsec:K-read} where, for a fixed reservoir channel $T$, we learn a readout $h \in \gHkp$ by minimizing the empirical risk $\hat{R}_N(H^T_h)$ under an RKHS norm constraint $\normkp{h} \leq \Lambda$, yielding an optimizer admitting a representer depending on the Gram matrix $K$ with elements $K_{ij} = \kappa\pth{H^T(\rvx_i), H^T(\rvx_j)}$ for windows $\rvx_i, \rvx_j \in \gD$. 

Since the conditions of the theorem satisfy those of Proposition~\ref{prop:set-injectivity}, the non-collision event $\mathrm{E_{NC}}$, which represents ``for any windows $\rvx \neq \rvx' \in \gD$, the feature vectors $m \circ H^T(\rvx)$ and $m \circ H^T(\rvx')$ do not collide'', happens with probability at least $1- \delta_{\pr}$. Condition on this event from now on.

As on $\mathrm{E_{NC}}$, the vectors $m \circ H^T(\rvx_i)$ are pairwise distinct and the Mat\'ern kernel is strictly positive definite on distinct inputs, the Gram matrix $K$ is positive definite, hence invertible. Let the label vector $\vy := (y_1, ..., y_N)^\top$, set $\valpha^\star := K^{-1} \vy$, and consider the representer
\begin{equation*}
    h^\star(\cdot) := \sum_{i=1}^N \alpha^\star_i \kappa\pth{H^T(\rvx_i), H^T(\cdot)} \in \gHkp.
\end{equation*}
Then we have $h^\star(H^T(\rvx_i)) = (K \valpha^\star)_i = y_i$ for all $i$, hence $\hat{R}_N(H^T_{h^\star}) = 0$. Additionally, we use $\normkp{h^\star}^2 = (\valpha^\star)^\top K \valpha^\star = \vy^\top K^{-1} \vy$ to define the threshold $\Lambda^\star := \normkp{h^\star}= \sqrt{\vy^\top K^{-1} \vy} > 0$.

Fix any $\Lambda \leq \Lambda^\star$ and let $h_\Lambda$ be the scaled function $h_\Lambda := \tfrac{\Lambda}{\Lambda^\star} h^\star$, which yields $\normkp{h_\Lambda} = \Lambda$. Using the squared loss $\ell(y,y') = (y-y')^2$ in~\cref{eq:empirical-risk-proxy} we get
\begin{equation*}
    \hat{R}_N\pth{H^T_{h_\Lambda}} = \frac{1}{N} \sum_{i=1}^N \pth{h_\Lambda(H^T(\rvx_i)) - y_i}^2 = \pth{1 - \frac{\Lambda}{\Lambda^\star}}^2 \frac{1}{N}\normeuc{\vy}^2.
\end{equation*}
Since $h_\Lambda$ is feasible for the norm-constrained problem, we get that
\begin{equation*}
    \min_{h\in \gHkp(\Lambda)} \hat{R}_N(H^T_h) \leq \pth{1 - \frac{\Lambda}{\Lambda^\star}}^2 \frac{1}{N}\normeuc{\vy}^2 = O\pth{\pth{1 - \frac{\Lambda}{\Lambda^\star}}^2},
\end{equation*}
since, as assumed in Section~\ref{subsec:learning-temporal-data}, the labels are bounded.

Hence, with probability at least $1-\delta_{\pr}$ (that of realization of the event $\mathrm{E_{NC}}$), the above inequality holds, thereby completing the proof.

\qed


\subsection{Weakly-dependent data}\label{subsec:weakly-dependent-data}
In this section, we provide a quick background on $\beta$-mixing processes and outline some results that we use in our proof of Theorem~\ref{thm:generalization}. 

\subsubsection{\texorpdfstring{$\beta$-}{β-}mixing processes}
Let $(\Omega, \gA, \sP)$ be a probability space. For two sub-$\sigma$-algebras $\gU, \gV \subset \gA$, the $\beta$-coefficient (absolute regularity coefficient)~\cite{Volkonskii_Rozanov_Yu_1959} is defined as 
\begin{equation}\label{eq:def-beta-mixing}
    \beta(\gU, \gV) := \frac{1}{2} \sup\acc{\sum_{i=1}^I \sum_{j=1}^J \abs{\sP(U_i \cap V_j) - \sP(U_i)\sP(V_j)}},
\end{equation}
where the supremum is taken over all finite measurable partitions $(U_i)_i$ of $\Omega$ from $\gU$ and $(V_j)_j$ from $\gV$. Equivalently, $\beta$ can be expressed in terms of a norm in total variation~\citep{dedecker2007} 
$$
\beta(\gU, \gV) = \norm{\sP_{\gU \otimes \gV} - \sP_{\gU} \otimes \sP_{\gV}}_{TV},
$$ 
where $\sP_{\gU}, \sP_{\gV}$ denote restrictions of $\sP$ to $\sigma$-fields $\gU, \gV$ and $\sP_{\gU \otimes \gV}$ is a law on the product $\sigma$-field defined on rectangles by $$\sP_{\gU \otimes \gV}(U,V) = \sP(U \cap V).$$

For a stationary random process $\rmX = (X_t)_{t \in \sZ_-}$, the mixing coefficients are obtained as
\begin{equation*}
    \beta_{\scriptscriptstyle \rmX}(k) = \sup_{t^\ast \in \sZ_-} \beta\Big(\sigma(X_t : t \leq t^\ast - k), \sigma(X_t : t \geq t^\ast) \Big),
\end{equation*}
where $\sigma(X_t : t)$ denotes the $\sigma$-field generated by the random process $(X_t : t)$. We say that $\rmX$ is $\beta$-mixing if $\beta_{\rmX}(k) \xrightarrow{k \to \infty} 0$, which means that past--future dependence (regularly) converges to $0$ as we increase the gap $k \in \sN$. Furthermore, $\rmX$ is called geometrically $\beta$-mixing if $\beta_{\rmX}(k) \leq \beta_0 e^{-\beta_1 k}$, and it is called algebraically $\beta$-mixing if $\beta_{\rmX}(k) \leq \beta_0 k^{-\beta_1}$ for some $\beta_0, \beta_1 > 0$.


\subsubsection{\texorpdfstring{$\beta$}{β}-mixing of the window process}\label{subsubsec:beta-mixing-window-process}
Let $w \in \sN$ be a window size, $s = w + g$ for some gap $g \in \sN$, and $\rmIO = ((X_t, Y_t): t \in \sZ_-)$ be a stationary $\beta$-mixing process. Define the windows process $\rmW := \pth{(\rmX_\tau^w, Y_\tau): \tau \in s \sZ_-}$, with $s\sZ_- := \{..., -2s, -s, 0\}$, which consists of dividing the process $\rmIO$ into $w$-sized windows
\begin{equation*}
    \rmX_\tau^w := \pth{X_{\tau - w + 1}, \ldots, X_\tau} \in \gW := (\gI)^w,
\end{equation*}
with indexing jumping by $s$.

The following result shows that the mixing coefficient $\beta_{\rmW}(k)$ of the windows process $\rmW$ is no larger than the mixing coefficient $\beta_{\rmIO}(ks-w)$ of the I/O process $\rmIO$, which we use later on.
\begin{claim}[$\beta$-mixing property of the windows process]\label{cl:beta-mixing-windows-process}
    Let $\rmIO = ((X_t, Y_t) : t \in \sZ_-)$ be the I/O process and $\rmW = \pth{(\rmX^w_\tau, Y_\tau) : \tau \in s \sZ_-}$, with $s = w + g$, for $g \in \sN^\ast$. Then we have
    \begin{equation}
        \beta_{\rmW}(k) \leq \beta_{\rmIO}(k s - w).
    \end{equation}
\end{claim}

\begin{proof}
Let $\rmW := (W_\tau: \tau \in s \sZ_-)$ where $W_\tau := (\rmX_\tau^w, Y_\tau)$. Let $\tau^\star \in s \sZ_-$ and define the past/future $\sigma-$fields for the windows process respectively as
\begin{equation*}
    \gF^- := \sigma\pth{W_\tau: \tau \leq \tau^\star - ks}, \qquad \gF^+ := \sigma\pth{W_\tau: \tau \geq \tau^\star}.
\end{equation*}
As each random window $W_\tau$ is measurable with respect to. $$\sigma\pth{IO_t : t \in \{\tau - w + 1, \ldots, \tau\}},$$ with elements $IO_t := (X_t, Y_t)$, we get 
\begin{equation*}
\begin{aligned}
    \gF^- \subset \sigma\pth{IO_t : t \leq \tau^\star - ks}, \qquad \gF^+ &\subset \sigma\pth{IO_t : t \geq \tau^\star - w + 1}\\
    &\subset \sigma\pth{IO_t : t \geq \tau^\star - w}.
\end{aligned}
\end{equation*}
We use monotonicity of $\beta(\cdot, \cdot)$ of~\cref{eq:def-beta-mixing} under increasing $\sigma-$fields (i.e., if $\gU \subset \gU'$ and $\gV \subset \gV'$, then $\beta(\gU, \gV) \leq \beta(\gU', \gV')$, to deduce that
\begin{equation*}
    \beta(\gF^-, \gF^+) \leq \beta\Big(\sigma\pth{IO_t : t \leq \ \tau^\star - ks}, \sigma\pth{IO_t : t \geq \tau^\star - w}\Big)
\end{equation*}
Setting $t^\star := \tau^\star - w$ we get $\tau^\star - ks = t^\star - (ks - w)$. Hence, we get
\begin{equation*}
\begin{aligned}
    \beta(\gF^-, \gF^+) &\leq \beta\Big(\sigma\pth{IO_t : t \leq \ t^\star - (ks - w)}, \sigma\pth{IO_t : t \geq t^\star}) \\
    &\leq \beta_{\rmIO}(ks - w),
\end{aligned}
\end{equation*}
where the last inequality holds by definition of the $\beta-$mixing coefficients of $\rmIO$. A supremum taken over $t^\star$ yields $\beta_{\rmW}(k) \leq \beta_{\rmIO}(ks - w)$ as claimed.
\end{proof}

\subsubsection{Non-i.i.d. Rademacher generalization bound}\label{subsubsec:non-iid-rademacher-generalization}
As our windows dataset is built from a single realization of the underlying I/O process, the samples $\{(\rvx_i, y_i)\}_{i=1}^N$ (equivalently $\{W_\tau\}$) are in general dependent. A tractable and standard way to account for such dependence is through the $\beta-$mixing coefficients of the window process $\rmW = (W_\tau : \tau \in s \sZ_-)$, as this mixing coefficient captures how quickly far-apart windows appear nearly independent. Non-i.i.d. Rademacher complexity bounds are specifically used for this context to characterize generalization under such setting. 
We provide below an adaptation of Theorem 2 of~\cite{mohri2008rademacher} to our window process $\rmW$.

\begin{theorem}[Mohri–Rostamizadeh~\cite{mohri2008rademacher}]\label{thm:mohri-rostamizadeh}
    Assume that $\rmW = (W_\tau : \tau \in s Z_-)$ is stationary and $\beta-$mixing with coefficients $(\beta_{\rmW}(k))_{k \geq 1}$. Let $\gC$ be a class of measurable functions $f: \gW \to [0, \Upsilon]$. Let $N = 2\mu$ be an even number with $\mu \in \sN^\star$. Define the total windows sample set $S_N$ and odd-indexed windows set $S_\mu \subset S_N$ resp. as
    \begin{equation*}
        S_N := \{\rvw_{-Ns}, \ldots \rvw_{-2s}, \rvw_{-s}\}, \quad S_{\mu} := \{\rvw_{-(2\mu - 1) s}, \ldots \rvw_{-3s}, \rvw_{-s}\}.
    \end{equation*}
    Let the statistical risk and empirical risk  of a hypothesis $f \in \gC$ be defined resp. as
    \begin{equation*}
        R(f) := \E[f(W_0)], \qquad \hat{R}_N(f) := \frac{1}{N} \sum_{i=1}^N f(\rvw_{-is}).
    \end{equation*}
    Let $\delta \in (0,1)$ be such that $\delta > 4 (\mu - 1) \beta_{\rmW}(1)$ and set $$\delta' := \delta - 4 (\mu - 1) \beta_{\rmW}(1).$$
    Then, with probability at least $1-\delta$ (over the draw of $S_N$), we have the following holding simultaneously for all $f \in \gC$:
    \begin{equation}\label{eq:Mohri-Rostamizadeh-risk-bound}
        R(f) \leq \hat{R}_N(f) + \Reb_{N/2}(\gC) + 3 \Upsilon \sqrt{\frac{\log(4/\delta')}{N}},
    \end{equation}
    where $\Reb_{N/2}(\gC)$ is the empirical Rademacher complexity on the subset $S_\mu$ defined as
    \begin{equation*}
        \Reb_{N/2}(\gC) := \E_{\epsilon}\cro{\sup_{f \in \gC} \; \frac{2}{\mu} \; \bigg| \sum_{\rvw \in S_\mu} \epsilon_i \, f(\rvw) \bigg|}, \qquad \epsilon_i \stackrel{\text{i.i.d.}}{\sim}\mathrm{Unif}\{\pm 1\}.
    \end{equation*}
\end{theorem}

\subsection{Proof of Theorem~\ref{thm:generalization}}
The proof of Theorem~\ref{thm:generalization} proceeds as follows.
\begin{enumerate}
    \item Bound the deviation between the window risk and the empirical risk, $\abs{R^w(H^T_h) - \hat{R}_N(H^T_h)}$, where $R^w(H^T_h)$ is the window risk given by
    \begin{equation}\label{eq:window-risk}
        R^w(H^T_h) := \E\cro{\ell\pth{H^T_h(\rmX^w), Y}},
    \end{equation}
    and where $(\rmX^w, Y)$ is a random pair consisting of an input window of size $w$ and its label.
    \item Bound the deviation between the true risk and the window risk, $\abs{R(H^T_h) - R^w(H^T_h)}$. 
    \item Finally, combine the two bounds via the triangle inequality to obtain
    \[
    \abs{R(H^T_h) - \hat{R}_N(H^T_h)}
    \leq
    \abs{R(H^T_h) - R^w(H^T_h)}
    +
    \abs{R^w(H^T_h) - \hat{R}_N(H^T_h)}.
    \]
\end{enumerate}

\subsubsection{Bounding the deviation of the window risk from the empirical one}
Below we show how this deviation behaves as a direct corollary of Theorem~\ref{thm:mohri-rostamizadeh}.

\begin{corollary}[Non-i.i.d.\ generalization for MSE and RKHS-ball readouts]
\label{cor:mse_rkhs_beta}
Let ${\rmIO}=((X_t,Y_t):t\in\sZ_-)$ be stationary and $\beta$-mixing, and let $s=w+g$ with $w,g\in \sN^\ast$.
Construct a strided windows dataset
$\gD=\{(\rvx_i,y_i)\}_{i=1}^{N}$ with $N=2\mu$ even by choosing indices
$t_1>t_2>\cdots>t_N$ such that $t_i-t_{i+1} = s$.
Fix a reservoir embedding $H^T$ and a kernel $\kappa$ as in~\cref{eq:matern-kernel},
and consider the RKHS ball
$\gHkp(\Lambda) = \{h\in\gHkp:\|h\|_\kappa\le \Lambda\}$.
Let $\ell$ be the squared loss $\ell(\hat y,y):=(\hat y-y)^2$.
Assume the labels are bounded $|Y_t|\le \Upsilon_Y$ and that the kernel is normalized so that
$\kappa(\rho,\rho)\le 1$ for all $\rho$ (respected by the Mat\'ern profile, $\kappa(\rho,\rho)=1$).
Let the window risk $R^w(H_h^T)$ be as defined in~\cref{eq:window-risk}.
If $\delta\in(0,1)$ satisfies
\[
\delta>4(\mu-1)\,\beta_{\rmIO}(g),
\qquad
\delta':=\delta-4(\mu-1)\,\beta_{\rmIO}(g),
\]
then with probability at least $1-\delta$, the following holds simultaneously for all
$h\in\gHkp(\Lambda)$:
\begin{equation}
    \abs{R^w(H_h^T)-\widehat R_N(H_h^T)}
    \;\le\; \Reb_{N/2}(\ell \circ \gHkp(\Lambda)) + M(N,w,g).
\end{equation}
where $\Reb_{N/2}(\ell \circ \gHkp(\Lambda))$ denotes the empirical Rademacher complexity of the loss-composed class $\ell \circ \gHkp(\Lambda)$ given by
\begin{equation}
    \Reb_{N/2}(\ell \circ \gHkp(\Lambda)) := \frac{4\Lambda(\Lambda+\Upsilon_Y)}{\sqrt{\mu}} = 4\sqrt{2}\,\frac{\Lambda(\Lambda+\Upsilon_Y)}{\sqrt{N}},
\end{equation}
and the mixing penalty $M(N,w,g)$ is given by
\begin{equation}
    M(N,w,g) := 3(\Lambda+\Upsilon_Y)^2\sqrt{\frac{\log(4/\delta')}{N}}.
\end{equation}
\end{corollary}

\begin{proof}

Let the hypothesis class $\gC := \{f_h : h \in \gHkp(\Lambda)\}$. By the reproducing property and the diagonal bound $\kappa(\rho, \rho) \leq 1$, we have for all $(\rvx,y) \in \gD$
\begin{equation*}
    \abs{h(H^T(\rvx))} \leq \normkp{h}\,\normkp{\kappa(H^T(\rvx), \cdot)} \leq \Lambda.
\end{equation*}
Since $|y| \leq \Upsilon_Y$, it follows that $0 \leq f_h(\rvx, y) \leq (\Lambda + \Upsilon_Y)^2$ for all $(\rvx, y)$, hence $\gC \subseteq [0,\Upsilon]^{\gW \times \sR}$ with $\Upsilon := (\Lambda + \Upsilon_Y)^2$.

Since $t_i-t_{i+1}=s$, the sequence $(\rvw_i)_{i=1}^N:=((\rvx_i,y_i))_{i=1}^N$ is a length-$N$ consecutive segment of the stationary window process $\rmW=(W_\tau:\tau\in s\sZ_-)$ (up to a deterministic time shift).
Therefore, Theorem~\ref{thm:mohri-rostamizadeh} applies directly to the class $\gC$ with mixing coefficient $\beta_{\rmW}(1)$.
Using Claim~\ref{cl:beta-mixing-windows-process}, we upper bound $\beta_{\rmW}(1)\le \beta_{\rmIO}(g)$,
hence $\delta'=\delta-4(\mu-1)\beta_{\rmIO}(g)$ and~\cref{eq:Mohri-Rostamizadeh-risk-bound} yields the bound
\begin{equation*}
    R^w(H_h^T) \leq \widehat R_N(H_h^T) + \Reb_{N/2}(\gC) + 3 \Upsilon \sqrt{\frac{\log(4/\delta')}{N}}.
\end{equation*}
It remains to upper bound the empirical Rademacher complexity $\Reb_{N/2}(\gC) =: \Reb_{N/2}(\ell \circ \gHkp(\Lambda))$. 

Let the odd-indexed sub-dataset be
$\gD^{\rm odd}:=\{(\rvx_{2i-1},y_{2i-1})\}_{i=1}^{\mu}$. By definition (Theorem~\ref{thm:mohri-rostamizadeh}), we have
\begin{equation*}
\begin{aligned}
    \Reb_{N/2}(\gC) = &\E_{\epsilon}\cro{\sup_{h \in \gHkp(\Lambda)} \; \frac{2}{\mu} \; \bigg| \sum_{i=1}^\mu \epsilon_i \, (h(\rho_{2i-1}) - y_{2i-1})^2 \bigg|}, \\
    &\epsilon_i \stackrel{\text{i.i.d.}}{\sim}\mathrm{Unif}\{\pm 1\},
\end{aligned}
\end{equation*}
where $\rho_{2i-1} := H^T(\rvx_{2i-1})$. Notice that for each $i$, the map $\psi_i(u) := (u - y_{2i-1})^2$ is $L$-Lipschitz on $[-\Lambda, \Lambda]$ with $L := 2 (\Lambda + \Upsilon_Y)$ as we have $\abs{\psi_i'(u)} = 2|u - y_{2i-1}| \leq 2 (\Lambda + \Upsilon_Y)$. By Talagrand's lemma (see, for example, Lemma~5.7 of~\cite{mohri2018foundations}), we have
\begin{equation*}
\begin{aligned}
    &\Reb_{N/2}(\gC) \leq 2 (\Lambda + \Upsilon_Y)\; \Reb_{N/2}(\gHkp(\Lambda)),\\
    \text{where} \;\; &\Reb_{N/2}(\gHkp(\Lambda)) :=\E_{\epsilon}\cro{\sup_{h \in \gHkp(\Lambda)} \; \frac{2}{\mu} \; \bigg| \sum_{i=1}^\mu \epsilon_i \, h(\rho_{2i-1}) \bigg|}.
\end{aligned}
\end{equation*}
$\Reb_{N/2}(\gHkp(\Lambda))$ is the empirical Rademacher complexity of the RKHS ball $\gHkp(\Lambda)$ which, by the reproducing property of the kernel $\kappa$ together with Cauchy--Schwarz, satisfies
\begin{equation*}
\begin{aligned}
    \Reb_{N/2}(\gHkp(\Lambda)) &\leq \frac{2 \Lambda}{\mu} \E_{\epsilon}\cro{\normkp{\sum_{i=1}^\mu \epsilon_i \kappa(\rho_{2i-1}, \cdot)}}\\
    & \leq \frac{2 \Lambda}{\mu} \sqrt{\E_\epsilon \normkp{\sum_{i=1}^\mu \epsilon_i \kappa(\rho_{2i-1}, \cdot)}^2}\\ 
    &= \frac{2 \Lambda}{\mu} \sqrt{\E_{\epsilon} \cro{\vepsilon^\top K \vepsilon}} = \frac{2 \Lambda}{\mu}\sqrt{\Trc K},
\end{aligned}
\end{equation*}
where $\vepsilon := (\epsilon_1, ..., \epsilon_\mu)^\top \in \{\pm 1\}^\mu$, $K$ is the Gram matrix $$K =(\kappa(\rho_{2i-1}, \rho_{2j-1}))_{i,j},$$ and the equality $\E_{\vepsilon}[\vepsilon^\top K \vepsilon] = \Trc K$ follows from independence and $\E[\epsilon_i\epsilon_j] = \delta_{ij}$. Again, as $\kappa$ is norm-diagonal, we have $\Trc K \leq \mu$, hence $\Reb_{N/2}(\gHkp(\Lambda)) \leq \tfrac{2\Lambda}{\sqrt{\mu}}$, which yields in turn
\begin{equation*}
    \Reb_{N/2}(\gC) \leq  \; \frac{4 \Lambda (\Lambda + \Upsilon_Y)}{\sqrt{\mu}}.
\end{equation*}

\end{proof}



\subsubsection{Bounding the deviation of the window risk from the statistical one}
Before stating a bound for such a deviation, we first show a geometric decaying of the dependence of the reservoir output on the far past, which is related to the exponentially fading memory property of our reservoir.
\begin{proposition}[Geometric decay of dependence of reservoir outputs on inputs]\label{prop:geometric-decay}
    Let $\rvx = (x_t : t\in \sZ_-) \in \gIZ$ be a time series, and let $\rvx^w : = (x_{-w + 1}, ..., x_0) \in (\gI)^w$ be its last truncated window. Let $H^T_h$ be a \textsc{QuaRK} reservoir as described in Section~\ref{sec:methodology} where $H^T$ is composed of $R \geq 1$ sub-reservoirs each of contraction factor $\lambda_r$, $\gO$ are the $k$-local observables measured per sub-reservoir, and $h \in \gHkp(\Lambda)$, where $\kappa$ is a Mat\'ern-based kernel as defined in~\cref{eq:matern-kernel} parameterized on $\nu > 1$, and $\xi > 0$. Then we have
    \begin{equation}
        \abs{H^T_h(\rvx) - H^T_h(\rvx^w)} \leq \frac{2 \Lambda}{\xi} \sqrt{\frac{\nu R |\gO|}{\nu - 1}} \lambda_\star^w,
    \end{equation}
    where $\lambda_\star := \max_r \lambda_r,$ and the reservoir is initialized for the finite window case to an arbitrary state $\rho'_w \in \gSH^R$.
\end{proposition}
\begin{proof}
    Let the sequence of states produced by the reservoir via consuming the full time series and the truncated window be respectively \begin{equation*} \begin{aligned} U^T(\rvx) &:= (\rho_t \in \gSH^R : t \in \sZ_-), \\ U^T(\rvx^w) &:= (\rho'_t \in \gSH^R : t \in \{-w+1, ..., 0\}), \end{aligned} \end{equation*} where, by definition, the truncated evolution $(U^T(x^w))_{-w}$ is initialized at some state $\rho'_{-w} \in \gSH^R$, and both $U^T(x)$ and $U^T(x^w)$ are then driven by the same input sequence over the last $w$ time steps.
    
    Since the composed reservoir $T(\rho, x)$ is $\lambda_\star-$contractive with respect to. the trace norm $\normTr{\cdot}$, with $\lambda^\star := \max_r \lambda_r$, we have 
    \begin{equation}\label{eq:back-recursion}
    \begin{aligned} \normTr{\rho_0 - \rho'_0} &\leq \lambda_\star \normTr{\rho_{-1} - \rho'_{-1}}\\ &\leq \lambda_\star^2 \normTr{\rho_{-2} - \rho'_{-2}}\\ &\;\;\vdots\\ &\leq \lambda_\star^w \normTr{\rho_{-w} - \rho'_{-w}} \leq  2 \lambda_\star^w, \end{aligned}
    \end{equation}
    where the last inequality leverages the fact that the diameter of the space of density operators is $2$ with respect to. the trace norm.

   We now account for the readout. By definition, we have
$$\abs{H^T_h(\rvx) - H^T_h(\rvx^w)}  = \abs{h(\rho_0) - h(\rho'_0)}.$$
Using the reproducing property in $\gHkp$ we get
\begin{equation}\label{eq:reprod-rho-0}
\begin{aligned}
    \abs{h(\rho_0) - h(\rho_0')} &= \abs{\Braket{h, \kappa(\rho_0, \cdot) - \kappa(\rho_0', \cdot)}_{\scriptscriptstyle \kappa}}\\
    &\leq \normkp{h} \normkp{\kappa(\rho_0, \cdot) - \kappa(\rho_0', \cdot)},
\end{aligned}
\end{equation}
with $\Braket{\cdot, \cdot}_{\scriptscriptstyle \kappa}$ being the RKHS inner product. Furthermore, we have
\begin{equation}\label{eq:kernel-to-matern-profile}
\begin{aligned}
    \normkp{\kappa(\rho_0, \cdot) - \kappa(\rho_0', \cdot)}^2 &= \kappa(\rho_0,\rho_0) + \kappa(\rho'_0,\rho'_0) - 2 \kappa(\rho_0,\rho'_0)\\
    &= 2 \pth{\varphi(0) - \varphi\pth{\normeuc{m(\rho_0) - m(\rho'_0)}}},
\end{aligned}
\end{equation}
where~\cref{eq:matern-kernel} was used. Define the constant
$$
L_\kappa := \sup_{s > 0} \frac{\sqrt{2 \pth{\varphi(0) - \varphi(s)}}}{s}.
$$
We show in the next paragraph that $L_\kappa \leq \sqrt{-\varphi''(0)},$ where $\varphi''(0)$ is the second derivative of $\varphi$ at the origin, which exists for the Mat\'ern profile $\varphi$ as $\nu > 1$ (i.e., $\varphi \in C^2$ at $0$).

To show that $L_\kappa \leq \sqrt{-\varphi''(0)}$, we use the standard Bochner~\cite{Hofmann_Schölkopf_Smola_2008} result stating that for stationary covariance (kernel) functions such as $\psi(\eta) := \varphi(\normeuc{\eta})$, there exists a finite nonnegative symmetric spectral measure $\gamma$ such that 
\begin{equation*}
    \psi(\eta) = \int_{\sR^d} e^{i \omega^\top \eta} d\gamma(\omega) = \int_{\sR^d} \cos(\omega^\top \eta) d\gamma(\omega),
\end{equation*}
where the last equality holds as $\psi$ is real and even. Hence we have
\begin{equation*}
    \varphi(0) - \varphi(\normeuc{\eta}) = \psi(0) - \psi(\eta) = \int_{\sR^d} \pth{1 - \cos(\omega^\top \eta)} d\gamma(\omega).
\end{equation*}
Using the elementary inequality $1- \cos t \leq t^2/2$ (valid for all $t$) yields
\begin{equation}\label{eq:inequality-matern-profile}
\begin{aligned}
    \varphi(0) - \varphi(\normeuc{\eta}) &\leq \frac{1}{2} \int_{\sR^d} (\omega^\top \eta)^2 d\gamma(\omega)\\
    &\leq \frac{\normeuc{\eta}^2}{2} \int_{\sR^d} (\omega^\top e)^2 d\gamma(\omega), \quad e := \tfrac{\eta}{\normeuc{\eta}}.
\end{aligned}
\end{equation}
Consider the 1D restriction $g(t) = \psi(t e) = \varphi(\abs{t})$. Differentiating the spectral representation twice at point $t = 0$ (legitimate as $\varphi''(0)$ exists), we get
\begin{equation*}
    g''(0) = \frac{d^2}{dt^2} \int \cos(t \omega^\top e) d\gamma(\omega) \bigg|_{t=0} = - \int (\omega^\top e)^2 d\gamma(\omega) = - \varphi''(0),
\end{equation*}
where the last equality holds as $g''(0) = \varphi''(0)$. The last inequality of~\cref{eq:inequality-matern-profile} therefore yields $2 (\varphi(0)-\varphi(s)) \leq -\varphi''(0) s^2$, which directly leads to
\begin{equation}\label{eq:L-kappa-bound}
    L_\kappa  \leq \sqrt{-\varphi''(0)}.
\end{equation}
Hence setting $s = \normeuc{m(\rho_0) - m(\rho'_0)}$ yields
\begin{equation*}
\begin{aligned}
    \sqrt{2 \pth{\varphi(0) - \varphi\pth{\normeuc{m(\rho_0) - m(\rho'_0)}}}} &\leq L_\kappa \normeuc{m(\rho_0) - m(\rho'_0)} \\
    &\leq \normeuc{m(\rho_0) - m(\rho'_0)} \sqrt{-\varphi''(0)}.
\end{aligned}
\end{equation*}

   
Furthermore, let $z := \sqrt{2\nu}\,s/\xi$ so that the Mat\'ern profile (defined in~\cref{eq:matern-profile}) becomes
\[
\varphi(s)=\frac{2^{1-\nu}}{\Gamma(\nu)}\, z^\nu K_\nu(z).
\]
For $\nu>1$, the small-$z$ expansion of $K_\nu$ gives
\[
z^\nu K_\nu(z)=2^{\nu-1}\Gamma(\nu)\left(1+\frac{z^2}{4(1-\nu)}+\mathcal{O}(z^4)\right)\qquad(z\to 0).
\]
Substituting into the definition of $\varphi$ yields
\[
\varphi(s)=1-\frac{1}{4(\nu-1)}z^2+\mathcal{O}(z^4)
=1-\frac{\nu}{2(\nu-1)\xi^2}s^2+\mathcal{O}(s^4).
\]
Hence $\varphi''(0)=-\frac{\nu}{(\nu-1)\xi^2}$, i.e.,
\[
-\varphi''(0)=\frac{1}{\xi^2}\frac{\nu}{\nu-1}.
\]
Therefore, from~\cref{eq:L-kappa-bound} we get $L_\kappa \leq \tfrac{1}{\xi} \sqrt{\tfrac{\nu}{\nu - 1}}$, and hence,
combining~\cref{eq:reprod-rho-0} and~\cref{eq:kernel-to-matern-profile}, we get
\begin{equation}\label{eq:readout-to-moments-bound}
\begin{aligned}
    \abs{h(\rho_0) - h(\rho_0')} &\leq \normkp{h} \normkp{\kappa(\rho_0, \cdot) - \kappa(\rho_0', \cdot)}\\
    &\leq \normkp{h} L_\kappa \normeuc{m(\rho_0) - m(\rho'_0)}\\
    &\leq \normkp{h} \frac{1}{\xi} \sqrt{\frac{\nu}{\nu - 1}} \normeuc{m(\rho_0) - m(\rho'_0)}.
\end{aligned}       
\end{equation}

Finally, for the moment map $m(\rho) = \pth{\Tr[O_\ell \rho]}_{\ell = 1}^{R |\gO|}$, where $R |\gO|$ is the number of observables measured across all sub-reservoirs, and where each $O_\ell$ is formed from a $k$-local observable (from $\gO$) padded by identities to act on the full space $\gSH^R$ (this does not change the operator norm), we have for each $\ell$ by H\"older's inequality for Schatten norms
\begin{equation}
    \abs{\Tr\cro{O_\ell (\rho_0 - \rho'_0)}} \leq \norm{O_\ell}_{\infty} \normTr{\rho_0 - \rho'_0}.
\end{equation}
As the operator norm of a Pauli string satisfies $\norm{O_\ell}_\infty = 1$, we get
\begin{equation}\label{eq:moments-to-states}
    \normeuc{m(\rho_0) - m(\rho'_0)} \!=\! \pth{\sum_{\ell=1}^{R |\gO|} \abs{\Tr\cro{O_\ell (\rho_0 - \rho'_0)}}^2 }^{1/2} \!\!\!\! \leq \sqrt{R |\gO|} \normTr{\rho_0 - \rho'_0}.
\end{equation}
Combining Eqs.~(\ref{eq:back-recursion}),~(\ref{eq:readout-to-moments-bound}), and~(\ref{eq:moments-to-states}), along with the fact that $\normkp{h} \leq \Lambda$, we get the claimed bound
\begin{equation*}
    \abs{h(\rho_0) - h(\rho_0')} \leq \frac{\Lambda}{\xi} \sqrt{\frac{\nu}{\nu - 1}} \cdot \sqrt{R |\gO|} \normTr{\rho_0 - \rho'_0} \leq \frac{2 \Lambda}{\xi} \sqrt{\frac{\nu R |\gO|}{\nu - 1}} \lambda_\star^w.
\end{equation*}
\end{proof}


We now provide the bound on the deviation $\abs{R(H^T_h) - R^w(H^T_h)}$.
\begin{corollary}[Deviation of the window risk from the statistical one]\label{cor:deviation-window-risk-true}
    Let the conditions of Proposition~\ref{prop:geometric-decay} apply, with the additional assumption that the considered time series is drawn from the stationary process $\rmIO = ((X_t,Y_t): t \in \sZ_-)$ and that the outputs satisfy $|Y_t| \leq \Upsilon_Y$. Then we have
    \begin{equation}
        \abs{R(H^T_h) - R^w(H^T_h)} \leq  \frac{4 \Lambda (\Lambda + \Upsilon_Y)}{\xi} \sqrt{\frac{\nu R |\gO|}{\nu - 1}} \lambda_\star^w.
    \end{equation}
\end{corollary}
\begin{proof}
    We have
    \begin{equation*}
        R(H^T_h) - R^w(H^T_h) = \E\cro{\ell\pth{H^T_h(\rmX), Y_0}} - \E\cro{\ell\pth{H^T_h(\rmX^w_0), Y_0}},
    \end{equation*}
    where by stationarity we write $\rmX = (\ldots, X_{-w-1}, X_{-w}, \rmX^w_0)$. Hence we get
    \begin{equation*}
    \begin{aligned}
        \abs{R(H^T_h) - R^w(H^T_h)} &= \abs{\E\cro{\ell\pth{H^T_h(\rmX), Y_0} - \ell\pth{H^T_h(\rmX^w_0), Y_0}}}\\
        & \leq \E\cro{\abs{\ell\pth{H^T_h(\rmX), Y_0} - \ell\pth{H^T_h(\rmX^w_0), Y_0}}}\\
        &\leq \E\cro{L \abs{H^T_h(\rmX) - H^T_h(\rmX^w_0)}}\\
        &\leq L \frac{2 \Lambda}{\xi} \sqrt{\frac{\nu R |\gO|}{\nu - 1}} \lambda_\star^w,
    \end{aligned}
    \end{equation*}
    where the last inequality uses that the squared loss is $L$-Lipschitz on $[-\Lambda,\Lambda]\times[-\Upsilon_Y,\Upsilon_Y]$, with $L := 2 (\Lambda + \Upsilon_Y)$ (see the third paragraph of the proof of Corollary~\ref{cor:mse_rkhs_beta}), together with Proposition~\ref{prop:geometric-decay}.
\end{proof}


We finally are now ready to provide proof of Theorem~\ref{thm:generalization}.

\subsubsection{Proof of Theorem~\ref{thm:generalization}}\label{subsubsec:proof-thm-generalization}
Since conditions of Theorem~\ref{thm:generalization} satisfy both conditions of Corollary~\ref{cor:mse_rkhs_beta} and Corollary~\ref{cor:deviation-window-risk-true}, and from 
$$\abs{R(H^T_h) - R_N(H^T_h)} \leq \abs{R(H^T_h) - R^w(H^T_h)} + \abs{R^w(H^T_h) - R_N(H^T_h)},$$
we conclude that 
\begin{equation*}
    \abs{R(H^T_h) - R_N(H^T_h)} \leq \Reb_{N/2}(\ell \circ \gHkp(\Lambda)) + M(N,w,g, \delta') + G(h, \Lambda, T)
\end{equation*}
where $\Reb_{N/2}(\ell \circ \gHkp(\Lambda))$ denotes the empirical Rademacher complexity of the loss-composed class $\ell \circ \gHkp(\Lambda)$ given by
$$\Reb_{N/2}(\ell \circ \gHkp(\Lambda)) := 4\sqrt{2}\,\frac{\Lambda(\Lambda+\Upsilon_Y)}{\sqrt{N}},$$
 the mixing penalty $M(N,w,g)$ given by
$$ M(N, w, g) := 3(\Lambda+\Upsilon_Y)^2\sqrt{\frac{\log(4/\delta')}{N}}$$
and the window truncation geometrically decaying penalty given by
$$\mathfrak{G}(\kappa, \Lambda, T) := \frac{4 \Lambda (\Lambda + \Upsilon_Y)}{\xi} \sqrt{\frac{\nu R |\gO|}{\nu - 1}} \lambda_\star^w.$$
\qed
